\chapter{Conclusion and Scope for Future Work}

\section{Introduction}

This chapter concludes the research paper by summarizing the contribution of the study, clarifying the scope boundaries within which the findings should be interpreted, and identifying directions for future extension. The goal is to bring the analytical, managerial, and technical strands of the project together in a balanced final view while preserving the practical value of the insights generated.

\section{Overall Contribution of the Study}

This study demonstrates that Google Play Store review analytics can provide meaningful insight into customer perception of Zomato. By combining sentiment analysis with theme-based interpretation, the project identifies where customer dissatisfaction is concentrated and where the brand retains relative strengths. Refund and cancellation, customer support, and delivery emerge as the most important areas requiring managerial attention, while app experience stands out as a comparatively positive aspect of the customer journey \cite{lemon2016,keller2013}.

The broader contribution of the study lies in showing that NLP-based review analysis is useful not only for technical classification, but also for strategic decision support. When properly structured, app review data can help convert large-scale customer feedback into actionable evidence for improving brand strategy, customer experience, and service performance \cite{liu2012,verhoef2010}. In this way, the project extends beyond model building into business interpretation and managerial usability.

\section{Study Boundaries}

The findings of the study should be interpreted in light of the chosen research frame. The analysis is centered on Google Play Store reviews, which provide a large and highly relevant source of app-based customer feedback. This focus is appropriate for the objectives of the project because the study is concerned with Zomato's digital customer voice in an app-mediated service setting.

Similarly, the sentiment labeling approach relies on review ratings as a structured proxy for overall sentiment, supported by textual cross-checks. This makes the dataset scalable and consistent for large-scale analysis. The theme classification method is also intentionally transparent, because it uses explicit keyword rules and a first-match-wins logic that can be explained and reproduced clearly in an MBA project context.

These choices do not weaken the study; rather, they define its scope. The project is designed as a practical, explainable, and decision-support-oriented review analytics study rather than as a purely experimental deep-learning investigation. Its value lies in interpretability, reproducibility, and direct business usefulness.

\section{Scope for Future Work}

The present study also opens several directions for future extension. First, the analysis can be expanded beyond Google Play Store reviews to include additional digital feedback channels such as App Store reviews, social media posts, customer complaint records, or support chat transcripts where appropriate and accessible.

Second, future work can explore more advanced multi-label theme classification so that a single review can be associated with more than one issue category when customers discuss overlapping concerns. Third, the sentiment framework can be extended with more sophisticated text-based models, including transformer-based approaches discussed in the broader NLP literature \cite{devlin2019,jurafsky2024}, if the research objective shifts toward deeper model experimentation.

Fourth, a future study may examine theme trends at a more granular level such as by geography, order type, campaign period, or service event. Fifth, the competitor comparison can be extended beyond one rival brand to build a broader market benchmark. These directions would allow the current project to evolve from a focused applied study into a more comprehensive customer intelligence framework.

\section{Final Conclusion}

The final conclusion of the study is that large-scale customer review data can be transformed into meaningful managerial insight when sentiment analysis, theme classification, benchmarking, and dashboard visualization are combined in a structured way. The Zomato case demonstrates that review analytics can reveal not only whether customers are satisfied or dissatisfied, but also why those perceptions are formed and which operational themes matter most.

For the present project, the strongest message is clear: the digital interface is not the central source of dissatisfaction. Instead, the major strategic pressure points lie in refund handling, support quality, and delivery execution. By acting on those areas while preserving its strengths in app experience, Zomato can use customer voice data not merely as feedback, but as a guide for more informed brand and service decisions.
